\ifdefined\gkver\else\def\gkver{student}\fi
\documentclass[\gkver]{gaokaozhenti}
\begin{document}

\chapter{2025年湖北}

\begin{enumerate}
\item
%% number: 1
%% paperName: 2025年湖北高考第 1 题 4 分
%% typeId: 1
%% type: 单选题
%% chapter: 原子和原子核
%% point: 半衰期
%% method: 无
%% score: 4
%% degree: 750
%% duplicateId: 0
%% body:
PET（正电子发射断层成像）是核医学科重要的影像学诊断工具，其检查原理是将含放射性同位素（如：\ce{^18_9 F}）的物质注入人体参与人体代谢，从而达到诊断的目的。
\ce{^18_9 F}的衰变方程为
\ce{^18_9 F -> X + ^0_1e + ^0_0 \nu}，
其中 \ce{^0_0 \nu}是中微子。
已知 \ce{^18_9 F}的半衰期是$110$分钟。下列说法正确的是 \xzanswer{C}

\fourchoices[answer=C]
{\ce{X}为 \ce{^17_8O}}
{该反应为核聚变反应}
{$1$克 \ce{^18_9 F}经$110$分钟剩下$0.5$克 \ce{^18_9 F}}
{该反应产生的 \ce{^0_0 \nu}在磁场中会发生偏转}

%% answer: C
%% memo:
\memoanswer{
	A．根据质量数与电荷数守恒可知，该物质为 \ce{^{18}_{8} O}，故 A 错误；
	
	B．核聚变是轻核结合成重核的过程（如氢弹原理）。本题中的衰变是单个原子核自发转变为另一种原子核，属于放射性衰变（具体为 $\beta^{+}$ 衰变），而非核聚变，故 B 错误；
	
	C. $ 1 \Ug $ 该物质经过 $ 110 \Umin $ 即一个衰变周期，则有一半发生衰变，该物质质量变为 $ 0.5 \Ug $，故 C 正确；
	
	D．\ce{^{0}_{0} \nu} 不带电，在磁场中不偏转，故 D 错误。
	
	故选 C。
}

\item
%% number: 2
%% paperName: 2025年湖北高考第 2 题 4 分
%% typeId: 1
%% type: 单选题
%% chapter: 曲线运动
%% point: 天体运动
%% method: 无
%% score: 4
%% degree: 650
%% duplicateId: 0
%% body:
甲、乙两行星绕某恒星做圆周运动，甲的轨道半径比乙的小。忽略两行星之间的万有引力作用，下列说法正确的是 \xzanswer{A} 

\fourchoices[answer=A]
{甲运动的周期比乙的小}
{甲运动的线速度比乙的小}
{甲运动的角速度比乙的小}
{甲运动的向心加速度比乙的小}


%% answer: A

%% memo:
\memoanswer{
	根据卫星做圆周运动的向心力等于万有引力可知

	$G\frac{Mm}{r^{2}}=m\frac{v^{2}}{r}=m\omega^{2}r=m\frac{4\pi^{2}}{T^{2}}r=ma$

	可得 $T = 2\pi \sqrt{\frac{r^{3}}{GM}}$，$v = \sqrt{\frac{GM}{r}}$，$\omega = \sqrt{\frac{GM}{r^{3}}}$，$a = \frac{GM}{r^{2}}$

	因 $r_\text{甲} < r_\text{乙}$，可知卫星甲、乙运动的周期 $T_\text{甲} < T_\text{乙}$，线速度关系 $v_\text{甲} > v_\text{乙}$，角速度关系 $\omega_\text{甲} > \omega_\text{乙}$，向心加速度关系 $a_\text{甲} < a_\text{乙}$。

	故选 A。
}

\item 
%% number: 3
%% paperName: 2025年湖北高考第 3 题 4 分
%% typeId: 1
%% type: 单选题
%% chapter: 气体性质
%% point: 热力学第一定律
%% method: 无
%% score: 4
%% degree: 650
%% duplicateId: 0
%% body:
如图所示，内壁光滑的汽缸内用活塞密封一定量理想气体，汽缸和活塞均绝热。用电热丝对密封气体加热，并在活塞上施加一外力$F$，使气体的热力学温度缓慢增大到初态的$2$倍，同时其体积缓慢减小。关于此过程，下列说法正确的是 \xzanswer{B} 


\onepicture[w=3.43cm, h=4.53cm]{figs/01.png}


\fourchoices[answer=B]
{外力$F$保持不变}
{密封气体内能增加}
{密封气体对外做正功}
{密封气体的末态压强是初态的$2$倍}


%% answer: B
%% memo:
\memoanswer{
	A．气体温度升高，体积减小，根据 $\frac{{pV}}{T} = C$ 可知气体压强变大，则外力 $F$ 增加，选项 A 错误；
	
	B．气体温度升高，则气体内能变大，即 $\Delta U$ 增加，选项 B 正确；
	
	C．气体体积减小，则外界对气体做功，选项 C 错误；
	
	D．根据 $ \frac{{pV}}{T} = C $，热力学温度变为原来的 $ 2 $ 倍，体积减小，则气体压强大于原来的 $ 2 $ 倍，选项 D 错误。
	
	故选 B。
}

\item 
%% number: 4
%% paperName: 2025年湖北高考第 4 题 4 分
%% typeId: 1
%% type: 单选题
%% chapter: 磁场
%% point: 磁场的叠加
%% method: 无
%% score: 4
%% degree: 650
%% duplicateId: 0
%% body:
如图所示，在磁感应强度大小为$ B $的匀强磁场中，放置一通电圆线圈，圆心为$ O $点，线圈平面与磁场垂直。在圆线圈的轴线上有$ M $和$ N $两点，它们到$ O $点的距离相等。已知$ M $点的总磁感应强度大小为零，则$ N $点的总磁感应强度大小为
\xzanswer{A} 

\onepicture[w=5.27cm, h=3.48cm]{figs/02.png}


\fourchoices[answer=A]
{$0$}
{$B$}
{$2B$}
{$3B$}


%% answer: A
%% memo:
\memoanswer{
	由右手螺旋定则及对称性可知，环形电流在 $ N $ 点产生的磁场，磁感应强度与 $ M $ 点等大同向。由于 $ M $ 点磁感应强度为零，由矢量合成法则可知环境中匀强磁场与 $ M $ 点磁场等大反向，即匀强磁场与 $ N $ 点的磁场等大反向，$ N $ 点的磁感应强度为$ 0 $。
	
	故选 A。
}
\memoanswer[shui]{
注意这里，站在$N$点或$M$点看环形电流，一个顺时针，一个逆时针，两者电流方向刚好相反，由于两者方向是相对的，因而在外看来是同向的。
}

\item 
%% number: 5
%% paperName: 2025年湖北高考第 5 题 4 分
%% typeId: 1
%% type: 单选题
%% chapter: 电磁感应
%% point: 电磁感应的电量
%% method: 无
%% score: 4
%% degree: 550
%% duplicateId: 0
%% body:
如图 \ref{2025湖北03a} 所示，相距$ L $的两足够长平行金属导轨放在同一水平面内，两长度均为$ L $、电阻均为$ R $的金属棒$ ab $、$ cd $垂直跨放在两导轨上，金属棒与导轨接触良好。导轨电阻忽略不计。导轨间存在与导轨平面垂直的匀强磁场，其磁感应强度大小$ B $随时间变化的图像如图 \ref{2025湖北03b} 所示，$ t=T $时刻，$ B=0 $。$ t=0 $时刻，两棒相距$ x_{0} $，$ ab $棒速度为零，$ cd $棒速度方向水平向右，并与棒垂直，则$ 0 \sim T $时间内流过回路的电荷量为 \xzanswer{B} 


\twopicture[labelA=2025湖北03a, labelB=2025湖北03b, hA=3.41cm, hB=4.11cm]
{figs/03a.png}{figs/03b.png}

\fourchoices[answer=B]
{$\frac{B_0Lx_0}{4R}$}
{$\frac{B_0Lx_0}{2R}$}
{$\frac{B_0Lx_0}R$}
{$\frac{2B_0Lx_0}R$}

%% answer: B
%% memo:
\memoanswer{
	通过导体的电荷量 $q = \overline{I} \Delta t$

	而 $\overline{I} = \frac{\overline{E}}{2R}$

	$t=T$ 时，磁感应强度为零，故 $\overline{E}=\frac{\Delta\Phi}{\Delta t}=\frac{B_{0}Lx_{0}-0}{T}$

	联立以上各式，可得 $q = \frac{B_{0}Lx_{0}}{2R}$

	故选 B。
}
\memoanswer[shui]{
	由$ \odif{q} =\frac{\odif{\Phi}}{R_\text{总}}$，得
$\Delta q =\frac{\Delta \Phi}{R_\text{总}}$，而$\Delta \Phi$只与初未磁通量有关，与中间的变化过程无关。
}

\item 
%% number: 6
%% paperName: 2025年湖北高考第 6 题 4 分
%% typeId: 1
%% type: 单选题
%% chapter: 曲线运动
%% point: 斜抛运动
%% method: 无
%% score: 4
%% degree: 450
%% duplicateId: 0
%% body:
某网球运动员两次击球时，击球点离网的水平距离均为$ L $，离地高度分别为$L/2$ 、$ L $，网球离开球拍瞬间的速度大小相等，方向分别斜向上、斜向下，且与水平方向夹角均为$ \theta $。击球后网球均刚好直接掠过球网，运动轨迹平面与球网垂直，忽略空气阻力，$ \tan \theta $的值为 \xzanswer{C} 


\fourchoices[answer=C]
{$\frac{1}{2}$}
{$\frac{1}{3}$}
{$\frac{1}{4}$}
{$\frac{1}{6}$}

%% answer: C

%% memo:
\memoanswer{
	网球水平方向上做匀速直线运动，有 $t=\frac{L}{v_{0}\cos\theta}$

	设球网高度为 $h$，则对 A 点发出的球，有

	$L-h=v_{0}\sin\theta\cdot t+\frac{1}{2}gt^{2}$

	对 B 点发出的球，有 $\frac{L}{2}-h=-v_{0}\sin\theta\cdot t+\frac{1}{2}gt^{2}$

	联立以上各式，可得 $\tan\theta=\frac{1}{4}$

	故选 C。
}

\item 
%% number: 7
%% paperName: 2025年湖北高考第 7 题 4 分
%% typeId: 1
%% type: 单选题
%% chapter: 运动和力
%% point: 加速减速模型
%% method: 无
%% score: 4
%% degree: 350
%% duplicateId: 0
%% body:
一个宽为$ L $的双轨推拉门由两扇宽为$L/2$的门板组成。门处于关闭状态，其俯视图如图 \ref{2025湖北04a} 所示。某同学用与门板平行的水平恒定拉力作用在一门板上，一段时间后撤去拉力，该门板完全运动到另一边，且恰好不与门框发生碰撞，其俯视图如图 \ref{2025湖北04b} 所示。门板在运动过程中受到的阻力与其重力大小之比为$ \mu $，重力加速度大小为$ g $。若要门板的整个运动过程用时尽量短，则所用时间趋近于 \xzanswer{B} 

\twopicture[labelA=2025湖北04a, labelB=2025湖北04b, w=5.80cm]
{figs/04a.png}{figs/04b.png}


\fourchoices[answer=B]
{$\sqrt{\frac{L}{2\mu g}}$}
{$\sqrt{\frac{L}{\mu g}}$}
{$\sqrt{\frac{2L}{\mu g}}$}
{$2\sqrt{\frac{L}{\mu g}}$}


%% answer: B
%% memo:
\memoanswer{
	设拉力为 $F$，作用时间为 $t_{1}$，撤去外力后 AB 运动的时间为 $t_{2}$，AB 运动过程的最大速度为 $v_{m}$，则由动量定理，有 $\left(F - \mu mg\right)t_{1} = mv_{m}$

	得 $t_{1}=\frac{mv_{m}}{F-\mu mg}$

	撤销拉力后，有 $\mu mgt_{2}=mv_{m}$

	得 $t_{2}=\frac{v_{m}}{\mu g}$

	对于全过程，有 $Ft_{1} = \mu mgt$

	得 $F = \frac{\mu mgt}{t_{1}}$

	对于全过程有 $\frac{L}{2} = \frac{v_{m}t}{2}$

	故 AB 运动的总时间

	$t = t_{1} + t_{2} = \frac{mv_{m}}{F - \mu mg} + \frac{v_{m}}{\mu g} = \frac{mv_{m}}{\frac{\mu mgt}{t_{1}} - \mu mg} + \frac{v_{m}}{\mu g} = \frac{v_{m}}{\mu g} \cdot \frac{t}{t - t_{1}} = \frac{v_{m}t}{\mu g} \cdot \frac{1}{t - t_{1}} = \frac{L}{\mu gt_{2}}$

	可知当 $t_{2}$ 越大时，$t$ 越小，当 $t_{2} = t$ 时，$t$ 取最小值。

	则 $t_{\min}^{2} = \frac{L}{\mu g}$

	则 $t_{\min} = \sqrt{\frac{L}{\mu g}}$

	故选 B。
}
\memoanswer[shui]{
可以使用函数法，但推荐图像法，面积固定找最短底边。
}

\item 
%% number: 8
%% paperName: 2025年湖北高考第 8 题 4 分
%% typeId: 2
%% type: 多选题
%% chapter: 交流电
%% point: 变压器
%% method: 无
%% score: 4
%% degree: 450
%% duplicateId: 0
%% body:
在如图所示的输电线路中，交流发电机的输出电压一定，两变压器均为理想变压器，左侧升压变压器的原副线圈匝数分别为$ n_{1} $和$ n_{2} $，两变压器间输电线路电阻为$ r $。下列说法正确的是 \xzanswer{AC} 

\onepicture[w=10.54cm, h=2.92cm]{figs/05.png}

\fourchoices[answer=AC]
{仅增加用户数，$ r $消耗的功率增大}
{仅增加用户数，用户端的电压增大}
{仅适当增加$ n_{2} $，用户端的电压增大}
{仅适当增加$ n_{2} $，整个电路消耗的电功率减小}


%% answer: AC


%% memo:
\memoanswer{
	A．整个电路物理量标注如图

	\onepicture[w=10.54cm]{figs/08_an.jpg}

	设用户端总电阻为 $R_\text{用户}$，降压器输入端等效为电阻 $R_{x}$，则有 $I_{3}^{2}R_{x}=I_{4}^{2}R_\text{用户}$

	因为 $I_{3}n_{3}=I_{4}n_{4}$

	其中 $I_{2}=I_{3}$，则对于输电线有

	$I_{2}=\frac{U_{2}}{r+R_{x}}$

	$I_{2}=\frac{U_{2}}{r+\left(\frac{n_{3}}{n_{4}}\right)^{2}R_\text{用户}}$

	联立整理可得

	仅增加用户数，则 $R_\text{用户}$ 减小，可知 $I_{2}$ 增大，根据 $\Delta P = I_{2}^{2} r$

	可知 $r$ 消耗的功率增大，故 A 正确；

	B．仅增加用户数，$I_{2}$ 增大，根据 $U_{3}=U_{2}-I_{2}r$

	可知 $U_{3}$ 减小，根据用户端的电压

	$U_{4}=\frac{n_{4}}{n_{3}}U_{3}$

	可知用户端的电压减小，故 B 错误；

	C．仅适当增加 $n_{2}$，根据

	$U_{2}=\frac{n_{2}}{n_{1}}U_{1}$

	$I_{2}=\frac{U_{2}}{r+\left(\frac{n_{3}}{n_{4}}\right)^{2}R_\text{用户}}$

	可知 $U_{2}$ 增大，根据

	可知 $I_{2}$ 增大，根据 $U_{3}=I_{2}R_{x}$

	$U_{4}=\frac{n_{4}}{n_{3}}U_{3}$

	可知 $U_{3}$ 增大，根据用户端的电压

	可知用户端的电压增大，故 C 正确；

	D．由 C 选项可知，仅适当增加 $n_{2}$，$I_{2}$ 增大，整个电路消耗的电功率 $P_{1}=P_{2}=U_{2}I_{2}$

	由于 $U_{2}$ 增大，故整个电路消耗的电功率变大，故 D 错误。

	故选 AC。
}

\item 
%% number: 9
%% paperName: 2025年湖北高考第 9 题 4 分
%% typeId: 2
%% type: 多选题
%% chapter: 机械振动
%% point: 简谐振动
%% method: 无
%% score: 4
%% degree: 350
%% duplicateId: 0
%% body:
质量均为$ m $的小球$ a $和$ b $由劲度系数为$ k $的轻质弹簧连接，小球$ a $由不可伸长的细线悬挂在$ O $点，系统处于静止状态，如图所示。将小球$ b $竖直下拉长度$ l $后由静止释放。重力加速度大小为$ g $，忽略空气阻力，弹簧始终在弹性限度内。释放小球$ b $后 \xzanswer{AD} 

\onepicture[w=3.16cm, h=4.88cm]{figs/06.png}


\fourchoices[answer=AD]
{小球$ a $可能会运动}
{若小球$ b $做简谐运动，则其振幅为$\frac l2$}
{当且仅当$l\leqslant\frac{mg}k$ 时，小球 b 才能始终做简谐运动}
{当且仅当$l\leqslant\frac{2mg}k$时，小球 b 才能始终做简谐运动}


%% answer: AD
%% memo:
\memoanswer{
	B.如果 $ A $ 球不动而 $ B $ 球单独振动则 $ B $ 球做简谐振动，简谐振动的平衡位置合力为零，即 $ B $ 球初始时刻位置，则可知 $ B $ 的振幅为 $ l $， B  错误；
	
	ACD．$A$球发生运动的临界条件为弹簧对 $ A $ 球向上的弹力大于 $ A $ 球的重力，则此时对 $ A $ 球有 $ kx_{0} = mg $
	
	对 $ B $ 球有此时加速度 $ kx_{0} + mg = ma $
	
	由简谐振动的对称性可得向下拉到最低点松手释放的加速度也为 $a$，则有 $kl = ma$
	
	解得 $l = \frac{{2mg}}{k}$
	
	即 $l \leq \frac{{2mg}}{k}$，否则 $A$ 球会发生运动，AD 正确，C 错误。
	
	故选 AD。
}

\item 
%% number: 10
%% paperName: 2025年湖北高考第 10 题 6 分
%% typeId: 2
%% type: 多选题
%% chapter: 电场
%% point: 电场叠加
%% method: 等效替代
%% score: 6
%% degree: 350
%% duplicateId: 0
%% body:
如图所示，在$ xOy $平面内有一以$ O $点为中心的正五边形，顶点到$ O $点的距离为$ R $。在正五边形的顶点上顺时针方向依次固定电荷量为$ q $、$ 2q $、$ 3q $、$ 4q $、$ 5q $的正点电荷，且电荷量为$ 3q $的电荷在$ y $轴正半轴上。静电力常量为$ k $，则$ O $点处的电场强度 \xzanswer{AD} 

\onepicture[w=3.95cm, h=3.56cm]{figs/07.png}


\fourchoices[answer=AD]
{方向沿$ x $轴负方向}
{方向与$ x $轴负方向成$ 18 \degree $夹角斜向下}
{大小为$\frac{2kq}{R^2} ( \cos 54 \degree + \cos 18 \degree ) $}
{大小为$\frac{2kq}{R^{2}} (2 \cos 54 \degree + \cos 18 \degree ) $}


%% answer: AD
%% memo:
\memoanswer{
	由题意可知，如图
	
	\onepicture[w=4.74cm, h=3.48cm]{figs/08.png}
	
	将五个点电荷等效成
	
	\onepicture[w=4.48cm, h=4.06cm]{figs/09.png}
	
	五个点电荷与 $O$ 点距离为 $R$，设 $E_{0}=\frac{kq}{R^{2}}$
	
	
	则 $ O $ 点场强大小为 $ E =2 \times 2  E_{0} \cos 54 {^\circ} +2 E_{0} \cos 18 {^\circ} $
	
	代入可得 $ E=\frac{2kq}{R^{2}}\left(2\cos54^{\circ}+\cos18^{\circ}\right) $
	
	
	方向沿 $ x $ 轴负方向；
	
	故选 AD。
}
\memoanswer[shui]{
这里化到上面的最简后有些同学还会继续分解，如左边是两$-q$，右边是两$+q$，下面是$-q$和$+q$，由此易得$E=\frac{2kq}{R^2}( 2\cos 36\Udu \cos 18\Udu + \cos 54\Udu )$，积化和差后与上面答案一样。
}
\item
%% number: 11
%% paperName: 2025年湖北高考第 11 题 8 分
%% typeId: 4
%% type: 实验题
%% chapter: 电路
%% point: 测电动势与内阻
%% method: 无
%% score: 8
%% degree: 450
%% duplicateId: 0
%% body:
某实验小组为测量一节干电池的电动势$ E $和内阻$ r $，设计了如图 \ref{2025湖北10a} 所示电路，所用器材如下：干电池、智能手机、电流传感器、定值电阻$ R_{0} $、电阻箱、开关、导线等。按电路图连接电路，将智能手机与电流传感器通过蓝牙无线连接，闭合开关$ S $，逐次改变电阻箱的阻值$ R $，用智能手机记录对应的电流传感器测得的电流$ I $。回答下列问题：
\twopicture[labelA=2025湖北10a, labelB=2025湖北10b, hA=3.8cm, hB=4.5cm]
{figs/10a.png}{figs/10b.png}

\begin{enumerate}
\item
$ R_{0} $在电路中起 \tkanswer{保护} （填“保护”或“分流”）作用。

\item 
$\frac{1}{I}$与$ E $、$ r $、$ R $、$ R_{0} $的关系式为$ \frac{1}{I}= $ \tkanswer{$\frac{R}{E}+\frac{R_0+r}{E}$} 。

\item 
根据记录数据作出$\frac{1}{I} - R $图像，如图 \ref{2025湖北10b} 所示。已知$ R_{0} =9.0 \UO $，可得$ E=$ \tkanswer{$1.47$} $ \UV[a] $（保留三位有效数字），$ r= $ \tkanswer{$1.3$} $\UO[a] $（保留两位有效数字）。


\item 
电流传感器的电阻对本实验干电池内阻的测量结果 \tkanswer{有}（填“有”或“无”）影响。

\end{enumerate}



%% answer:
\jdanswer{
\begin{enumerate}
	\item 
	保护
	
	\item
	$\frac{R}{E} + \frac{{{R_0} + r}}{E}$
	
	\item
	$1.47$，$1.3$
	
	\item
	有
\end{enumerate}
}

%% memo:
\memoanswer{
\begin{enumerate}
	\item
	$R_{0}$ 与电阻箱串联，可知，$R_{0}$ 在电路中起保护作用。

	\item
	根据闭合电路欧姆定律 $E = I(R + R_{0} + r)$

	化简可得 $\frac{1}{I}=\frac{R}{E}+\frac{R_{0}+r}{E}$

	\item
	结合上述有 $\frac{1}{I}=\frac{R}{E}+\frac{R_{0}+r}{E}$

	结合图（b）有 $\frac{1}{E} = \frac{24 - 7}{25 - 0}\UV^{-1}$，$\frac{R_{0} + r}{E} = 7 \UA^{-1}$

	解得 $E \approx 1.47 \UV$，$r \approx 1.3 \UO$

	\item
	当电流传感器有内阻时，所测的电源内阻 $r_\text{测} = r_\text{真} + r_\text{传}$，导致电源内阻测量值偏大，即电流传感器的电阻对本实验干电池内阻的测量结果有影响。
\end{enumerate}
}

\item 
%% number: 12
%% paperName: 2025年湖北高考第 12 题 9 分
%% typeId: 4
%% type: 实验题
%% chapter: 功和能
%% point: 整体机械能守恒
%% method: 无
%% score: 9
%% degree: 450
%% duplicateId: 0
%% body:
某同学利用如图 \ref{2025湖北11} 所示的实验装置来测量重力加速度大小$ g $。细绳跨过固定在铁架台上不可转动的小圆柱体，两端各悬挂一个重锤。实验步骤如下：
\twopicture[labelA=2025湖北11, labelB=2025湖北12, hA=4cm, wB=6.5cm]
{figs/11.png}{figs/12.png}


①用游标卡尺测量遮光片的宽度$ d $。

②将遮光片固定在重锤$ 1 $上，用天平测量重锤$ 1 $和遮光片的总质量$ m $、重锤$ 2 $的质量$ M $（$ M>m $）。

③将光电门安装在铁架台上，将重锤$ 1 $压在桌面上，保持系统静止，重锤$ 2 $离地面足够高。用刻度尺测量遮光片中心到光电门的竖直距离$ H $。

④启动光电门，释放重锤$ 1 $，用毫秒计测出遮光片经过光电门所用时间$ t $。

⑤根据上述数据求出重力加速度$ g $。

⑥多次改变光电门高度，重复步骤，求出$ g $的平均值。


回答下列问题：
\begin{enumerate}
\item
测量$ d $时，游标卡尺的示数如图 \ref{2025湖北12} 所示，可知 \tkanswer{0.515} $ \Ucm[a] $。

\item 
重锤$ 1 $通过光电门时的速度大小为$ v= $ \tkanswer{$\frac{d}{t}$} （用遮光片$ d $、$ t $表示）。若不计摩擦，$ g $与$ m $、$ M $、$ d $、$ t $、$ H $的关系式为$g=$ \tkanswer{$\frac{(M+m)d^2}{2(M-m)Ht^2}$} 。

\item 
实验发现，当$ M $和$ m $之比接近于$ 1 $时，$ g $的测量值明显小于真实值。主要原因是圆柱体表面不光滑，导致跨过圆柱体的绳两端拉力不相等。理论分析表明，圆柱体与绳之间的动摩擦因数很小时，跨过圆柱体的绳两端拉力差$ \Delta T=4 \gamma  \frac{Mm}{M+m} g$，其中$ \gamma $是只与圆柱体表面动摩擦因数有关的常数。保持$ M+m=2 m_{0} $不变，其中$ M=(1+ \beta ) m_{0} $，$ m=(1 - \beta ) m_{0} $。$ \beta $足够小时，重锤运动的加速度大小可近似表示为$ a=( \beta - \gamma )g $。调整两重锤的质量，测得不同$ \beta $时重锤的加速度大小$ a $，结果如下表。根据表格数据，采用逐差法得到重力加速度大小$ g=$ \tkanswer{$9.81$} $\Umsq[a] $（保留三位有效数字）。


\begin{table}[htbp]
\centering 
\begin{tblr}{|c|c|c|c|c|}
\hline 
$\beta$ & 0.04 & 0.06 & 0.08 & 0.10
 \\
\hline
$a$ ($\Umsq[a]$) & 0.084 & 0.281 & 0.477 & 0.673\\ 
\hline 
\end{tblr}
\end{table} 


\end{enumerate}

%% answer:
\jdanswer{
\begin{enumerate}
	\item 
	$0.515$
	
	\item
	$\frac{d}{t}$，$\frac{{(M + m){d^2}}}{{2(M - m)H{t^2}}}$
	
	\item
	$9.81$
\end{enumerate}
}
%% memo:
\memoanswer{
\begin{enumerate}
	\item 
	根据游标卡尺的读数规律，该游标卡尺的读数为 $ 5 \Umm +0.05 \times 3 \Umm =5.15 \Umm =0.515 \Ucm $。
	
	\item
	根据光电门的测速原理，重锤 $1$ 通过光电门时的速度大小为 $v = \frac{d}{t}$。
	
	对重锤 $1$ 与重锤 $2$ 构成的系统进行分析，根据系统机械能守恒定律有($M - m$)$gH = \frac{1}{2}$($M + m$)$v^{2}$，其中 $v = \frac{d}{t}$。解得 $g = \frac{{(M + m){d^2}}}{{2(M - m)H{t^2}}}$。
	
	\item
	由于 $\gamma$ 是只与圆柱体表面动摩擦因数有关的常数，且有$a$ = ($\beta - \gamma$)$g = \beta g -\gamma g$。取表格从左至右四组数据分别为$a_{1}$、$a_{2}$、$a_{3}$、$a_{4}$和对应的$\beta_{1}$、$\beta_{2}$、$\beta_{3}$、$\beta_{4}$，利用表格中的数据，根据逐差法有 $a_{4} + a_{3} - a_{2} - a_{1}$ = ($\beta_{4} + \beta_{3} - \beta_{2} - \beta_{1}$)$g$，带入数据可则重力加速度 $g = 9.81 \Umsq$。
\end{enumerate}
}

\item 
%% number: 13
%% paperName: 2025年湖北高考第 13 题 9 分
%% typeId: 6
%% type: 计算题
%% chapter: 光学
%% point: 全反射
%% method: 无
%% score: 9
%% degree: 450
%% duplicateId: 0
%% body:
如图所示，三角形$ ABC $是三棱镜的横截面，$ AC=BC $，$ \angle C=30 \degree $，三棱镜放在平面镜上，$ AC $边紧贴镜面。在纸面内，一光线入射到镜面$ O $点，入射角为$ \alpha $，$ O $点离$ A $点足够近。已知三棱镜的折射率为$\sqrt{2}$。
\begin{enumerate}
	\item
	当$ \alpha =45 \degree $时，求光线从$ AB $边射入棱镜时折射角的正弦值；
	
	\item 
	若光线从$ AB $边折射后直接到达$ BC $边，并在$ BC $边刚好发生全反射，求此时的$ \alpha $值。
	
\end{enumerate}
\onepicture[h=2.9cm, align=right]{figs/13.png}

%% answer:
\jdanswer{
\begin{enumerate}
	\item
	$\sqrt{6}/4$
	\item 
	$60 \degree $
\end{enumerate}
}
%% memo:
\memoanswer{
	\onepicture[w=7.91cm, h=3.14cm]{figs/14.png}
}



\item 
%% number: 14
%% paperName: 2025年湖北高考第 14 题 16 分
%% typeId: 6
%% type: 计算题
%% chapter: 磁场
%% point: 带电粒子在磁场中的运动
%% method: 无
%% score: 16
%% degree: 350
%% duplicateId: 0
%% body:
如图所示，两平行虚线$ MN $、$ PQ $间无磁场。$ MN $左侧区域和$ PQ $右侧区域内均有垂直于纸面向外的匀强磁场，磁感应强度大小分别为$ B $和$ 2B $。一质量为$ m $、电荷量为$ q $的带正电粒子从$ MN $左侧$ O $点以大小为$ v_{0} $的初速度射出，方向平行于$ MN $向上。已知$ O $点到$ MN $的距离为$ \frac{3 m v_{0}} {2 q B} $，粒子能回到$ O $点，并在纸面内做周期性运动。不计重力，求：
\begin{enumerate}
	\item
	粒子在$ MN $左侧区域中运动轨迹的半径；
	\item 
	粒子第一次和第二次经过$ PQ $时位置的间距；
	\item 
	粒子的运动周期。
\end{enumerate}
\onepicture[h=4.3cm, align=right]{figs/15.png}

%% answer:
\jdanswer{
	\begin{enumerate}
		\item
		$R=\frac{mv_0}{qB} $
		\item 
		$ x=\frac{\sqrt{3}mv_0}{2qB}$
		\item 
		$ \frac{5\pi m}{3qB}+\frac{\sqrt{3}m}{qB}$
	\end{enumerate}
}



%% memo:
\memoanswer{
\begin{enumerate}
	\item
	粒子在左侧磁场中运动，根据洛伦兹力提供向心力有 $qv_{0}B=\frac{mv_{0}^{2}}{R}$

	可得 $R = \frac{mv_{0}}{qB}$

	\item
	粒子在左侧磁场运动，设从 $MN$ 射出时速度方向与 $MN$ 的夹角为 $\theta$，由于 $O$ 到 $MN$ 的距离 $d = \frac{3mv_{0}}{2qB}$，结合 $R = \frac{mv_{0}}{qB}$，根据几何关系可知 $\theta = 60^\circ$；

	粒子在 $MN$ 和 $PQ$ 之间做匀速直线运动，所以粒子从 $PQ$ 进入右侧磁场时与 $PQ$ 的夹角 $\theta = 60^\circ$；

	粒子在右侧磁场做匀速圆周运动有 $qv_{0} \cdot 2B = \frac{mv_{0}^{2}}{R'}$

	解得 $R' = \frac{mv_{0}}{2qB}$

	根据几何关系可知粒子第一次和第二次经过 $PQ$ 时位置的间距 $x = 3 R' = \frac{3mv_{0}}{2qB}$。

	\item
	由图可知粒子在左边磁场运动的时间

	$t_{1}=\frac{2}{3}T_{1}=\frac{2}{3}\times\frac{2\pi m}{qB}=\frac{4\pi m}{3qB}$

	粒子在右边磁场运动的时间

	$t_{2}=\frac{1}{3}T_{2}=\frac{1}{3}\times\frac{2\pi m}{2qB}=\frac{\pi m}{3qB}$

	根据对称性可知粒子在 $MN$ 左侧进出磁场的距离

	$x_{0}=\sqrt{3}R=\frac{3mv_{0}}{2qB}$

	所以粒子从 $MN$ 到 $PQ$ 过程中运动的距离为

	$l=\frac{\sqrt{3}R-\sqrt{3}R'}{2\cos\theta}=\frac{\sqrt{3}mv_{0}}{2qB}$

	\onepicture[h=4.3cm]{figs/14_an.jpg}

	粒子在 $MN$ 和 $PQ$ 之间运动的时间 $t_{3}=\frac{2l}{v_{0}}=\frac{\sqrt{3}m}{qB}$

	综上可知粒子完成完整运动回到 $O$ 点的周期为 $T = t_{1} + t_{2} + t_{3} = \frac{5\pi m}{3qB} + \frac{\sqrt{3}m}{qB}$。
\end{enumerate}
}

\item 
%% number: 15
%% paperName: 2025年湖北高考第 15 题 18 分
%% typeId: 6
%% type: 计算题
%% chapter: 运动和力
%% point: 动量和能量
%% method: 无
%% score: 18
%% degree: 250
%% duplicateId: 0
%% body:
如图所示，一足够长的平直木板放置在水平地面上，木板上有$ 3n $（$ n $是大于$ 1 $的正整数）个质量均为$ m $的相同小滑块，从左向右依次编号为$ 1 $、$ 2 $、$ \cdots $、$ 3n $，木板的质量为$ nm $。相邻滑块间的距离均为$ L $，木板与地面之间的动摩擦因数为$ \mu $，滑块与木板间的动摩擦因数为$ 2 \mu $。初始时木板和所有滑块均处于静止状态。现给第$ 1 $个滑块一个水平向右的初速度，大小为$\sqrt{\beta \mu gL }$（$ \beta $为足够大常数，$ g $为重力加速度大小）。滑块间的每次碰撞时间极短，碰后滑块均会粘在一起继续运动。最大静摩擦力等于滑动摩擦力。
\begin{enumerate}
	\item
	求第$ 1 $个滑块与第$ 2 $个滑块碰撞前瞬间，第$ 1 $个滑块的速度大小;
	
	\item 
	记木板滑动前第$ j $个滑块开始滑动时的速度为$ v_j $，第$ j+1 $个滑块开始滑动时的速度为$ v_{j+1} $。用已知量和$ v_j $表示$ v_{j+1} $。
	
	\item 
	若木板开始滑动后，滑块间恰好不再相碰，求$ \beta $的值。
	
	（参考公式：$ 1^2+2^2+ \cdots +k^2= \frac{1}{6} k(k+1)(2k+1) $
	
	
\end{enumerate}
\onepicture[h=1cm, align=right]{figs/16.png}

%% answer:
\jdanswer*{
	\begin{enumerate}
		\item
		$v_{1}=$ $\sqrt {( \beta - 4) \mu gL}$
		\item 
		$v_{j+ 1}= \frac j{j+ 1}\sqrt {v_{j}^{2}- 4\mu gL}$
		\item 
		$\beta=\frac{8n}{2n-1}(2n+1)^2 + \frac{4}{3}n(2n+1)(4n+1)=\frac{4n(2n+1)(8n^{2}+10n+5)}{3(2n-1)}$
	\end{enumerate}
}

%% memo:
\memoanswer{
\begin{enumerate}
	\item
	滑块 1 运动时，对木板的摩擦力为 $f_{1}=2\mu mg$

	地面对木板的摩擦力为 $f_{2}=4\mu nmg$

	所以此过程中木板保持不动；每个滑块之间距离为 $L$，所以对滑块 1 根据动能定理有

	$-2\mu mgL=\frac{1}{2}mv_{1}^{2}-\frac{1}{2}mv_{0}^{2}$

	解得 $v_{1} = \sqrt{(\beta - 4)\mu gL}$

	\item
	滑块间碰撞时间极短，碰后滑块粘在一起运动，若长木板不动，第 $j$ 个滑块开始运动时加速度为 $a_{j}=\frac{2\mu\cdot jmg}{jm}=2\mu g$

	根据运动学公式，第 $j$ 个滑块开始滑动到和第 $j+1$ 个滑块碰撞时，有 $v_{j}^{2}-v_{j}^{\prime 2}=2a_{j}L$

	第 $j$ 个滑块和第 $j+1$ 个滑块碰撞过程中动量守恒有 $jmv_{j}^{\prime}=(j+1)mv_{j+1}$

	联立可得 $v_{j+1}=\frac{j}{j+1}\sqrt{v_{j}^{2}-4\mu gL}$

	\item
	当第 $k$ 个木块开始滑动时，木板恰好要滑动，此时有 $2\mu \cdot kmg = \mu(nm + 3nm)g$

	解得 $k = 2n$（$n$ 为整数）

	则第 $k+1$ 个（即 $2n+1$）木块开始滑动时，木板开始滑动，要刚好不发生下一次碰撞，假设木板和剩下的木块不发生相对滑动，则

	$2\mu(2n+1)mg - \mu(nm + 3nm)g = \left[nm + (n-1)m\right]a_\text{木}$

	则 $a=\frac{2\mu g}{2n-1}<2\mu g$

	木板和剩下的木块不发生相对滑动。

	$a_{k+1}=\frac{2\mu\cdot\left(k+1\right)mg}{\left(k+1\right)m}=2\mu g$

	对前面 $k+1$ 个（即 $2n+1$）木块，有

	木板开始滑动时，刚好不发生下一次碰撞，则对前面 $k+1$ 个木块和 $k+2$ 个木块共速，且相对位移恰好为 $L$，则

	$\frac{v_{k+1}^{2}-v_\text{共}^{2}}{2a_{k+1}}-\frac{v_\text{共}^{2}}{2a}=L$

	则 $v_{k+1}^{2} = 2nv_\text{共}^{2} + 4\mu gL$

	又 $v_\text{共} = at = v_{k+1} - 2\mu gt$

	则 $v_\text{共} = \frac{v_{k+1}}{2n}$

	则 $v_{k+1}^{2} = \frac{8n\mu gL}{2n-1}$

	$j = 1$ 时，第一个滑块开始运动的速度 $v_{0} = \sqrt{\beta\mu gL}$，则 $v_{0}^{2} = \beta\mu gL$

	$j = 2$ 时，根据动量守恒定律可得 $mv_{1} = 2mv_{2}$

	可得第 2 个滑块开始运动的速度 $v_{2} = \frac{1}{2}\sqrt{(\beta - 4)\mu gL}$

	则 $4v_{2}^{2} = \beta\mu gL - 4\mu gL$

	由第二问可得，$(j+1)^{2} v_{j+1}^{2} = j^{2} v_{j}^{2} - 4j^{2}\mu gL$，则对第 3 个滑块到第 $k+1$ 个滑块有

	$9v_{3}^{2} = 4 \times v_{2}^{2} - 4 \times 4\mu gL$

	$16v_{4}^{2} = 9v_{3}^{2} - 9 \times 4\mu gL$

	$25v_{5}^{2} = 16 \times v_{4}^{2} - 16 \times 4\mu gL$

	……

	$(k+1)^{2} v_{k+1}^{2} = k^{2} v_{k}^{2} - k^{2} \times 4\mu gL$

	将从 $j = 2$ 到 $j = k+1$ 相关方程累积求和可得

	$(k+1)^{2} v_{k+1}^{2} = \beta\mu gL - 4\mu gL \left[1^{2} + 2^{2} + 3^{2} + 4^{2} + \cdots + k^{2}\right] = \beta\mu gL - \frac{2k(k+1)(2k+1)}{3}\mu gL$

	联立 $v_{k+1}^{2} = \frac{8n\mu gL}{2n-1}$，$k = 2n$

	可得 $\beta = \frac{4n(2n+1)(8n^{2} + 10n+5)}{3(2n-1)}$
\end{enumerate}
}
\memoanswer[shui]{
	整个题目过程比较清晰，计算有点繁琐，最后一问以木板为参考系，考虑惯性力，可以简化计算。
}

\end{enumerate}

\end{document}
